\subsection{Analysis Procedure}
\label{ssec:unf_strat}

The Pb-going FCal transverse energy is constructed by summing the transverse
projection of the electromagnetic-scale tower energies in the hemisphere facing
the Pb beam. This definition is consistent with previous ATLAS $p$+Pb
centrality measurements. In typical events containing a hard scattering, the
energy observed in the Pb-going ZDC is dominated by spectator neutrons
evaporating from the struck nucleus.

The ZDC energy scale is determined run by run from the response to a single
spectator neutron. The low-energy spectrum contains distinguishable neutron
peaks, and a simultaneous fit to the first three peaks and a background term
validates the calibration. The fitted single-neutron peak is constrained to the
nominal per-nucleon beam energy. The detector response below $1~\mathrm{TeV}$
is limited by detector noise and low-energy backgrounds; this region is treated
as a dedicated zero-energy bin. A residual in-time pileup contribution to the
\EZDC distribution is removed with a data-driven template method.

Finite jet-resolution effects cause migration in $x_p$. These effects are
corrected with a two-dimensional iterative Bayesian unfolding procedure using
RooUnfold. The response matrix contains $x_p$ and, in turn, either \EZDC or
\ETPB. Since the available simulation does not model nuclear breakup in the ZDC
or the full FCal response for the overlaid minimum-bias event, the calorimeter
observable is propagated through the unfolding with diagonal migration
matrices. This retains the measured correlation of calorimeter energy with
$x_p$ while correcting the jet-related kinematic migration.

The response matrix is reweighted as a function of reconstructed $x_p$ to
match the data spectrum. For the ZDC measurement, pileup-subtracted data
distributions are sampled to assign a ZDC energy to each simulated dijet event.
The procedure is repeated with independent random seeds, and closure tests
verify that the unfolding reproduces the particle-level simulation. Two
iterations are used, balancing the convergence of the unfolded result against
the growth of statistical uncertainty. Statistical uncertainties are evaluated
by bootstrap resampling of both the data and response matrix.
